\chapter{Experiments}
\label{chap:experiments}

\chapterabstract{This chapter defines the experimental protocol used to validate the thesis claims and reports the resulting quantitative and qualitative findings.}

\section{Research Questions}
\placeholder{State the questions the experiments answer. Each question should map to a table, figure, ablation, or qualitative analysis.}

\section{Datasets}
\placeholder{Describe all datasets, splits, annotations, preprocessing steps, and any generated data.}

\section{Baselines}
\placeholder{List baseline methods and explain how they were configured. Make comparisons fair and reproducible.}

\section{Metrics}
\placeholder{Define quantitative metrics and explain what each metric captures. Include limitations of each metric where relevant.}

\section{Experimental Setup}
\placeholder{Specify training schedule, optimizer, learning rate, batch size, number of runs, random seeds, hardware, and software versions.}

\begin{table}[t]
  \centering
  \caption[Experiment configuration.]{Example experiment configuration table. Replace with the final setup.}
  \label{tab:experiment-config}
  \begin{tabular}{@{}ll@{}}
    \toprule
    Item & Value \\
    \midrule
    Dataset split & \placeholder{train / validation / test} \\
    Batch size & \placeholder{--} \\
    Optimizer & \placeholder{--} \\
    Hardware & \placeholder{--} \\
    \bottomrule
  \end{tabular}
\end{table}

\section{Main Results}
\label{sec:main-results}
\placeholder{Present the primary quantitative results first. Explain the trend, not just the numbers.}

\begin{table}[t]
  \centering
  \caption[Main results.]{Main result table placeholder.}
  \label{tab:main-results}
  \begin{tabular}{@{}lccc@{}}
    \toprule
    Method & Metric 1 & Metric 2 & Metric 3 \\
    \midrule
    Baseline & -- & -- & -- \\
    Proposed & -- & -- & -- \\
    \bottomrule
  \end{tabular}
\end{table}

\section{Ablation Studies}
\placeholder{Isolate the effect of each contribution. Keep ablations tied to the claims made in the introduction.}

\section{Qualitative Results}
\placeholder{Show representative successes and failures. Qualitative figures should make the method's behavior inspectable.}

\section{Limitations}
\placeholder{Discuss failure cases, data bias, compute cost, assumptions, and evaluation gaps. Strong limitations make the thesis more credible.}

\section{Discussion}
\placeholder{Connect the results back to the motivation. Explain what the experiments imply for the broader research problem.}
